\documentclass[varwidth=20cm, border=10pt]{standalone}
\usepackage{graphicx,capt-of}
\usepackage{multicol}

\newcounter{subfloat}
\renewcommand{\thesubfloat}{\alph{subfloat}}
\newcommand{\image}[2]{%
  \stepcounter{subfloat}%
  \begin{tabular}[t]{@{}c@{}}
  #2 \\
  (\thesubfloat) #1
  \end{tabular}%
}

\begin{document}

\begin{tabular}{cc}
\image{Plaintext image, 2000 by 1400 pixels, 24 bit color depth.}{%
  \includegraphics[width=.4\textwidth]{./Plaintext.png}
} & \image{ECB mode ciphertext, 5 pixel (120 bit) block size.}{%
  \includegraphics[width=.4\textwidth]{./Ciphertext5.png}
}\\
\image{ECB mode ciphertext, 30 pixel (720 bit) block size.}{%
  \includegraphics[width=.4\textwidth]{./Ciphertext30.png}
} & \image{ECB mode ciphertext, 100 pixel (2400 bit) block size.}{%
  \includegraphics[width=.4\textwidth]{./Ciphertext100.png}
} \\
\image{ECB mode ciphertext, 400 pixel (9600 bit) block size.}{%
  \includegraphics[width=.4\textwidth]{./Ciphertext400.png}
} &
\image{Ciphertext under idealized encryption.}{%
  \includegraphics[width=.4\textwidth]{./Random.png}
}\\
\end{tabular}

\end{document}
